GCC has a new infrastructure to support a link time optimization (LTO). The
infrastructure is designed to allow linking of large applications
using a special mode (WHOPR) which support parallelization of the compilation process.
In this paper we present overview of the design and
implementation of WHOPR and present test results of its behavior when optimizing large applications. We
give numbers on compile time, memory usage and code quality comparisons to the
classical file by file based optimization model. In particular we focus on Firefox web
browser. We show main problems seen only when compiling a large
application, such as startup time and code size growth.
